% R1 — Abstract
\begin{abstract}
Learning useful descriptions of movement from limited video data could broaden access to quantitative behavioral analysis. We evaluate a compact Skeleton Joint Embedding Predictive Architecture (S-JEPA) on gait videos labeled Normal, multiple sclerosis, or Parkinson's disease. In each split, the encoder learns from 24 source videos, comprising 45--58 clips, by predicting hidden skeleton features without condition labels. A separate linear classifier uses its frozen features to predict collection labels for recordings excluded from training. Evaluation across five source-grouped folds covers 88 usable clips from 41 videos. S-JEPA achieves source-weighted macro-F1 of 0.457, compared with 0.325 for a Random Forest using handcrafted joint-angle features and 0.452 for mean pose. With equal weight per clip, S-JEPA also exceeds the Random Forest, scoring 0.397 versus 0.321. These results show that a compact predictive skeleton model trained on a small video collection can produce useful features for held-out label prediction. Mean pose remains a strong comparison, and further controls are needed to isolate the contribution of pretraining and temporal information. The study supports representation learning for markerless behavioral analysis, with clinical interpretation requiring recordings with verified diagnoses and participant identities.
\end{abstract}
